\section{Background and Motivation}

This section should define the requirements engineering problem addressed by the paper and situate it within software engineering research. It should explain the relevant lifecycle stage, artifact types, stakeholders, and practical failure modes.

\subsection{Requirements Engineering Context}

Discuss the task in relation to established requirements engineering practices, standards, or artifacts such as stakeholder goals, user requirements, system requirements, software requirements specifications, traceability links, quality attributes, and validation criteria \cite{iso29148,nuseibeh2000requirements,wiegers2013software}.

\subsection{Motivating Example}

Include a compact motivating example that can be reused throughout the paper. The example should expose the problem, show why naive automation is insufficient, and preview how the proposed method improves the outcome.

\subsection{Research Gap}

State why existing tools, single-LLM workflows, rule-only methods, or retrieval-only methods are insufficient for the target task. Connect the gap to the proposed knowledge-driven or multi-agent design.

